\documentclass[10pt,a4paper,twoside,openright,titlepage,fleqn,%
               headinclude,,footinclude,BCOR5mm,%
               numbers=noenddot,cleardoublepage=empty,%
               tablecaptionabove]{scrbook}
\usepackage[dutch,polutonikogreek,british]{babel}
\usepackage[utf8]{inputenc}
\usepackage[LGRx,T1]{fontenc}
\usepackage[final]{microtype}
\usepackage{amsmath,amssymb,amsthm}
\usepackage{varioref}
\usepackage[style=philosophy-modern,hyperref,backref,square,natbib,ibidtracker=false]{biblatex}
\usepackage[tight,british]{minitoc}
\usepackage{wrapfig}
\usepackage{chngpage}
\usepackage{calc}
\usepackage{mflogo}
\usepackage{caption,listings,graphicx,subfig}
\usepackage{multicol}
\usepackage{makeidx}
\usepackage{xspace}
\usepackage{mparhack}
\usepackage{fixltx2e}
\usepackage{relsize}
\usepackage{lipsum}
\usepackage[eulerchapternumbers,subfig,beramono,eulermath,pdfspacing,listings]{classicthesis}
\usepackage{guit}
\usepackage{arsclassica}
\usepackage{graphicx}
\usepackage{tikz}
\usepackage{setspace}

% ********************************************************************
% classicthesis-preamble
% ********************************************************************

\newcommand{\myName}{Xavier Goas Aguililla\xspace}
\newcommand{\myTitle}{A corpus-based approach to the language of the papyri\xspace}
%\newcommand{\mySubTitle}{Methodological questions and select topics\xspace}
\newcommand{\myLocation}{Leuven\xspace}
\newcommand{\myUni}{KU Leuven\xspace}
\newcommand{\myFac}{Faculty of Arts\xspace}
\newcommand{\myDegree}{Master of Arts in de taal- en letterkunde\xspace}
\newcommand{\myProf}{dr. T. Van Hal\xspace}
\newcommand{\myDepartment}{Department of Greek Studies\xspace}

\newcommand{\myTime}{June 2012\xspace}

% ********************************************************************
% hyperref
% ******************************************************************** 
\hypersetup{%
    colorlinks=true, linktocpage=true, pdfstartpage=1, pdfstartview=FitV,%
    breaklinks=true, pdfpagemode=UseNone, pageanchor=true, pdfpagemode=UseOutlines,%
    plainpages=false, bookmarksnumbered, bookmarksopen=true, bookmarksopenlevel=1,%
    hypertexnames=true, pdfhighlight=/O,%
    urlcolor=webbrown, linkcolor=RoyalBlue, citecolor=RoyalBlue, pagecolor=RoyalBlue,%
% uncomment the following line if you want to have black links (e.g., for printing)
% urlcolor=Black, linkcolor=Black, citecolor=Black, pagecolor=Black,%
    pdftitle={\myTitle},%
    pdfauthor={\textcopyright\ \myName},%
    pdfsubject={},%
    pdfkeywords={},%
    pdfcreator={pdfLaTeX},%
    pdfproducer={LaTeX con hyperref e ClassicThesis}%
}

\hypersetup{citecolor=webgreen}
\hypersetup{hyperfootnotes=false,pdfpagelabels}


\newcommand{\mail}[1]{\href{mailto:#1}{\texttt{#1}}}



% ********************************************************************
% caption
% ********************************************************************
\captionsetup{format=hang,font=small}
\captionsetup[table]{skip=\medskipamount} 


% ********************************************************************
% makeidx, multicol
% ********************************************************************
\let\orgtheindex\theindex
\let\orgendtheindex\endtheindex
\def\theindex{%
	\def\twocolumn{\begin{multicols}{2}}%
	\def\onecolumn{}%
	\clearpage
	\orgtheindex
}
\def\endtheindex{%
	\end{multicols}%
	\orgendtheindex
}

\makeindex

% ********************************************************************
% listings
% ********************************************************************

\definecolor{lightergray}{gray}{0.99}

\lstset{language=[LaTeX]Tex,
    keywordstyle=\color{RoyalBlue},
    basicstyle=\normalfont\ttfamily,
    commentstyle=\color{Emerald}\ttfamily,
    stringstyle=\rmfamily,
    numbers=none,
    numberstyle=\scriptsize,
    stepnumber=5,
    numbersep=8pt,
    showstringspaces=false,
    breaklines=true,
    frameround=ftff,
    frame=lines,
    backgroundcolor=\color{lightergray}
}

\lstset{morekeywords=%
        {RequirePackage,newboolean,DeclareOption,setboolean,%
        ProcessOptions,PackageError,ifthenelse,boolean,%
        chapterNumber,sodef,textls,allcapsspacing,%
        MakeTextLowercase,orgtheindex,endtheindex,%
        @ifpackageloaded,undefined,sfdefault,%
        DeclareRobustCommand,spacedallcaps,%
        microtypesetup,MakeTextUppercase,lowsmallcapsspacing,%
        lowsmallcapsspacing,spacedlowsmallcaps,
        spacedlowsmallcaps,lehead,headmark,color,%
        headfont,partname,thepart,titleformat,part,
        titlerule,chapter,thechapter,thesection,%
        subsection,thesubsection,thesubsubsection,%
        paragraph,theparagraph,descriptionlabel,titlespacing,%
        graffito,lineskiplimit,finalhyphendemerits,%
        colorbox,captionsetup,labelitemi,%
        myincludegraphics,hypersetup,setlength,%
        definecolor,lsstyle,textssc,subsubsection,%
        graffito@setup,includegraphics,ifdefined,%
        myTitle,textcopyright,myName,lstset,lstnewenvironment,%
        setkeys,lst@BeginAlsoWriteFile,contentsname,%
        toc@heading,@ppljLaTeX,z@,check@mathfonts,%
        sf@size,ptctitle,mtctitle,stctitle,lst@intname,%
        @empty,math@fontsfalse,@ppljscTeX,@iwonaTeX,%
        @iwonascLaTeX,@ctTeX,tw@,ct@sc,@ctTeX,f@family,%
        f@shape,ct@sc,ctLaTeX,ctLaTeXe,@twoe,@sctwoe,%
        texorpdfstring,m@th,ctTeX,@mkboth,ProvidesPackage,%
        theindex,PackageInfo,PackageWarningNoLine,%
        mtifont,mtcindent,@iwonaLaTeX,@ppljTeX,@iwonascTeX,%
        rohead,orgendtheindex,@ppljscLaTeX,%
        @ifclassloaded,toc@headingbkORrp,backreftwosep,%
        backrefalt,backreflastsep,areaset,pnumfont,%
        arsincludegraphics,ExecuteOptions,PackageWarning,%
        MessageBreak,ars@@includegraphics,ifcld@backref,rofoot,formatchapter},
        commentstyle=\color{Emerald}\ttfamily,%
        frame=lines}

\lstset{basicstyle=\normalfont\ttfamily}
\lstset{flexiblecolumns=true}
\lstset{moredelim={[is][\normalfont\itshape]{/*}{*/}}}

\DeclareRobustCommand*{\pacchetto}[1]{{\normalfont\ttfamily#1}%
\index{Pacchetto!#1@\texttt{#1}}%
\index{#1@\texttt{#1}}}


\DeclareRobustCommand*{\classicthesis}{ClassicThesis}


\DeclareRobustCommand*{\bibtex}{\textsc{Bib}\TeX%
\index{bibtex@\textsc{Bib}\protect\TeX}%
}

\DeclareRobustCommand*{\amseuler}{\protect\AmS{} Euler%
\index{AmS Euler@\protect\AmS~Euler}%
\index{Font!AmS Euler@\protect\AmS~Euler}}

\lstset{basicstyle=\normalfont\ttfamily}
\lstset{flexiblecolumns=false}
\lstset{moredelim={[is][\ttfamily]{!?}{?!}}} 
\lstset{escapeinside={?*}{*?}}
\lstset{firstnumber=last}

\lstnewenvironment{code}% 
{\setkeys{lst}{columns=fullflexible,keepspaces=true}%
\lstset{basicstyle=\small\ttfamily}%
}{}

\lstset{extendedchars} 
\lstnewenvironment{sidebyside}{% 
    \global\let\lst@intname\@empty 
    \setbox\z@=\hbox\bgroup 
    \setkeys{lst}{columns=fullflexible,% 
    linewidth=0.45\linewidth,keepspaces=true,%
    breaklines=true,% 
    breakindent=0pt,%
    boxpos=t,%
    basicstyle=\small\ttfamily
}% 
    \lst@BeginAlsoWriteFile{\jobname.tmp}% 
}{% 
    \lst@EndWriteFile\egroup 
        \begin{center}% 
            \begin{minipage}{0.45\linewidth}% 
                \hbox to\linewidth{\box\z@\hss} 
            \end{minipage}% 
            \qquad 
            \begin{minipage}{0.45\linewidth}%
            \setkeys{lst}{frame=none}% 
                \leavevmode \catcode`\^^M=5\relax 
                \small\input{\jobname.tmp}% 
            \end{minipage}% 
        \end{center}% 
} 


\newcommand{\omissis}{[\dots\negthinspace]}

\graphicspath{{images/}}

\newcommand{\meta}[1]{$\langle${\normalfont\itshape#1}$\rangle$}
\lstset{escapeinside={?*}{*?}}

\lstset{moredelim={[is][\ttfamily]{!?}{?!}}}

\DeclareRobustCommand*{\miktex}{MiK\TeX%
\index{miktex@MiK\protect\TeX}%
}

\DeclareRobustCommand*{\metafont}{\MF%
\index{METAFONT@\protect\MF}%
}

\DeclareRobustCommand*{\metapost}{\MP%
\index{METAPOST@\protect\MP}%
}

\DeclareRobustCommand*{\texlive}{\TeX{}~Live%
\index{texlive@\protect\TeX{}~Live}%
}


%%%%%%%%%%%%
% biblatex %
%%%%%%%%%%%%

\bibliography{Bibliography}

\renewcommand{\nameyeardelim}{, }

\defbibheading{bibliography}{%
\cleardoublepage
\manualmark
\phantomsection
\addcontentsline{toc}{chapter}{\tocEntry{\bibname}}
\chapter*{\bibname\markboth{\spacedlowsmallcaps{\bibname}}
{\spacedlowsmallcaps{\bibname}}}}     

  \DeclareCiteCommand{\citeyearpar}[\mkbibparens] 
  {\boolfalse{citetracker}% 
   \boolfalse{pagetracker}% 
   \usebibmacro{prenote}} 
  {\printtext[bibhyperref]{\printfield{year}}} 
  {\multicitedelim} 
  {\usebibmacro{postnote}} 

\makeatletter 
  \DeclareCiteCommand{\citetalias} 
  {\usebibmacro{prenote}} 
  {\usebibmacro{citeindex}% 
   \bibhyperref{\@citealias{\thefield{entrykey}}}} 
  {\multicitedelim} 
  {\usebibmacro{postnote}} 
\makeatother 


%%%%%%%%%%%%%%%%%%
% other commands %
%%%%%%%%%%%%%%%%%%

\newcommand{\ita}[1]{% 
  \begin{otherlanguage*}{italian}#1\end{otherlanguage*}}
  
\DeclareRobustCommand*{\pkgname}[1]{{\normalfont\sffamily#1}%
\index{Package!#1@\textsf{#1}}%
\index{#1@\textsf{#1}}}

\DeclareRobustCommand*{\envname}[1]{{\normalfont\ttfamily#1}%
\index{Environment!#1@\texttt{#1}}%
\index{#1@\texttt{#1}}}

\DeclareRobustCommand*{\optname}[1]{{\normalfont\ttfamily#1}%
\index{Option!#1@\texttt{#1}}%
\index{#1@\texttt{#1}}}

\DeclareRobustCommand*{\clsname}[1]{{\normalfont\sffamily#1}%
\index{Class!#1@\textsf{#1}}%
\index{#1@\textsf{#1}}}

\DeclareRobustCommand*{\cmdname}[1]{\mbox{\lstinline!\\#1!}%
\index{#1@\texttt{\hspace*{-1.2ex}\textbackslash#1}}}

\DeclareRobustCommand*{\classicthesis}{Classic\-Thesis}

\DeclareRobustCommand*{\arsclassica}{{\normalfont\sffamily ArsClassica}}

\DeclareRobustCommand*{\miktex}{MiK\TeX%
\index{miktex@MiK\protect\TeX}}

\DeclareRobustCommand*{\texlive}{\TeX{}~Live%
\index{texlive@\protect\TeX{}~Live}}





\newcommand{\Greek}{\fontencoding{LGR}\selectfont}
\newcommand{\Latin}{\fontencoding{T1}\selectfont}
\newcommand\fullsc[1]{\scalebox{1.06}[1.09]{\textsc{#1}}}

\begin{document}
\pagestyle{plain}
\dominitoc

%******************************************************************
% front matter
%******************************************************************
\frontmatter
\begin{titlepage}
\changetext{}{}{}{((\paperwidth  - \textwidth) / 2) - \oddsidemargin - \hoffset - 1in}{}
\null\vfill
\begin{center}
\large
\sffamily

\vspace{2cm}
{\Large\spacedlowsmallcaps{\myName}} \\
\bigskip


{\huge\spacedlowsmallcaps{\myTitle} \\
}

\bigskip

{\Large\spacedlowsmallcaps{\mySubTitle}} \\

\bigskip

    
\vspace{7cm}

\begin{tabular} {cc}
\parbox{0.3\textwidth}{\includegraphics[width=3cm]{sedes}}
&
\parbox{0.7\textwidth}{\normalsize{Master's thesis submitted in 
	partial fulfilment of the requirements for the degree 
		of}\\ \\ {\Large\spacedlowsmallcaps{\myDegree}}\\ 

					{\normalsize
					Supervisor: \myProf \\
					\myUni \\
					\myFac \\
					\myDepartment \\
					}}
\end{tabular}
\vfill
\sc{Leuven, 2013}
\end{center}
\end{titlepage}




%*******************************************************
% Titleback
%*******************************************************
\thispagestyle{empty}

\hfill
\vfill

\noindent\myName:
\textit{\myTitle,} %\mySubTitle,
\textcopyright\ \myTime. 

\noindent Typeset with pdfTeX 3.1415926-2.4-1.40.13 using a modified version of Lorenzo Pantieri's \href{http://www.ctan.org/tex-archive/macros/latex/contrib/arsclassica/}{ArsClassica package}.

\medskip
\noindent{\spacedlowsmallcaps{E-mail}}: \\
\mail{xavier.goas@student.kuleuven.be}

\bigskip

\clearpage
%*******************************************************
% Abstract+Sommario
%*******************************************************

\begingroup
\let\clearpage\relax
\let\cleardoublepage\relax
\let\cleardoublepage\relax

% \chapter*{Abstract}

% \vfill

\selectlanguage{dutch}
\pdfbookmark[1]{Samenvatting}{Samenvatting}
\chapter*{Samenvatting}

%motivation
Ondanks de bewezen diensten van corpusgebaseerde methoden in het
taalkundig onderzoek is er een groot gebrek aan geannoteerde digitale
corpora voor het Oudgrieks. Grote hoeveelheden tekst zijn ondertussen
digitaal beschikbaar, maar door gebrek aan mankracht is het niet
mogelijk deze handmatig te analyseren. Een betrouwbare methode om dit
probleem automatisch aan te pakken, al zou deze niet perfect zijn, zou
een stap in de goede richting zijn. Een belangrijk pijnpunt is dat dit
gebrek aan geannoteerde corpora het ook zeer moeilijk maakt om dit
soort methode te ontwikkelen.

% problem statement & approach
In dit werk trachten we dit pijnpunt te omzeilen door een dergelijke
methode te ontwikkelen met behulp van spitstechnieken uit de
artificiële intelligentie en bij wijze van \textit{case study} toe te
passen op een corpus gedigitaliseerde documentaire papyri. Deze
methode berust op een implementatie van een zgn. artificieel neuraal
netwerk, d.i. een wiskundig model dat in staat is om patronen te
herkennen en te leren. We passen de methode beschreven in
\cite{collobert-2011} toe op het Oudgrieks.

Het leerproces van dit netwerk wordt opgedeeld in twee fasen. In een
eerste wordt gebruik gemaakt van een groot ongeannoteerd corpus om via
een eenvoudig criterium een wiskundig model op te bouwen dat woorden
kan kaderen binnen het taalgebruik in het geheel.  Dit model kent
kansen toe aan reeksen woorden al naargelang die al dan niet 'correct'
zijn door te schatten of gelijkaardige reeksen kunnen voorkomen in het
corpus.

Een tweede fase schakelt over op een kleiner, geannoteerd corpus. Na
het defini\"eren van een wiskundige voorstelling voor morfologische en
syntactische categorie\"en worden meerdere netwerken ge\"initialiseerd
die gebruik maken van de voorheen opgebouwde wiskundige
representatie. De modellen worden afgestemd op de woord-annotatieparen
en be\"invloeden elkaar onderling om zo goed mogelijk gebruik te maken
van patronen die van potentieel nut zijn voor de taak van elk netwerk.

We waren van plan om het zo bekomen model toe te passen op een
digitaal beschikbaar corpus papyri, en dit na afloop in een
verdeelbaar formaat om te zetten en te opensourcen voor verdere
verwerking. De finale resultaten waren echter teleurstellend: het
systeem is te eenvoudig om correcte morfologische en syntactische
inferenties te maken en lijdt nog steeds onder de kleine hoeveelheden
geannoteerde tekst die beschikbaar zijn voor Oudgrieks. De evolutie
van het model dat ontwikkeld werd in de eerste fase van het leerproces
is echter positief en stemt hoopvol. Met langere rekentijd kan dit
type model zeker van nut zijn in toekomstig onderzoek. We stellen ook
een aantal pistes voor om ook de tweede fase van het leerproces te
verbeteren.






% Problem statement:
% What problem are you trying to solve? What is the scope of your work (a generalized approach, or for a specific situation)? Be careful not to use too much jargon. In some cases it is appropriate to put the problem statement before the motivation, but usually this only works if most readers already understand why the problem is important.
% Motivation:
% Why do we care about the problem and the results? If the problem isn't obviously "interesting" it might be better to put motivation first; but if your work is incremental progress on a problem that is widely recognized as important, then it is probably better to put the problem statement first to indicate which piece of the larger problem you are breaking off to work on. This section should include the importance of your work, the difficulty of the area, and the impact it might have if successful.
% Approach:
% How did you go about solving or making progress on the problem? Did you use simulation, analytic models, prototype construction, or analysis of field data for an actual product? What was the extent of your work (did you look at one application program or a hundred programs in twenty different programming languages?) What important variables did you control, ignore, or measure?
% Results:
% What's the answer? Specifically, most good computer architecture papers conclude that something is so many percent faster, cheaper, smaller, or otherwise better than something else. Put the result there, in numbers. Avoid vague, hand-waving results such as "very", "small", or "significant." If you must be vague, you are only given license to do so when you can talk about orders-of-magnitude improvement. There is a tension here in that you should not provide numbers that can be easily misinterpreted, but on the other hand you don't have room for all the caveats.
% Conclusions:
% What are the implications of your answer? Is it going to change the world (unlikely), be a significant "win", be a nice hack, or simply serve as a road sign indicating that this path is a waste of time (all of the previous results are useful). Are your results general, potentially generalizable, or specific to a particular case?

\selectlanguage{british}

\endgroup			

\vfill
Deze masterproef bevat \textbf{59840
} tekens.


\pagestyle{scrheadings}
%************************************************
\chapter{Preface}
\label{chap:preface}
\mtcaddchapter
%************************************************

The papyri are an invaluable resource for documenting the history and 
evolution of the Greek language. The recently published \emph{The 
Language of the Papyri} [\cite{lpapyri}] has placed the 
spotlight firmly on the potential of this field while at the 
same time pointing out the regrettable lack of recent scholarly 
work available in it.

Though there is now a collection of texts available that is 
well-formatted and can easily be converted into an adequate 
corpus for linguistic research, comparatively few scholars are 
interested in exploiting the available resources; and none, to 
my knowledge, have attempted to do so.


A first necessity is, of course, a broad \emph{status quaestionis} yet 
in a less traditional sense than one might expect: the focus here lies 
more upon technical and technological concerns than specialised 
monographs for which data was gathered manually. Nevertheless, I have 
for the sake of completeness chosen to include a comprehensive 
bibliography of the field for clear reference. The term bibliography 
would be rather less appropriate in describing the array of databases 
and linguistic tools available to us - it is rather a select catalogue 
documenting the most relevant items in the instrumentarium in some 
detail while providing a briefly annotated list of other useful 
resources.

A second chapter is dedicated

\begin{enumerate}
\item a survey of previous studies on the language of the Greek papyri 
and an analysis of the methodology used therein;
\item a critical evaluation of the methods used by aforementioned 
studies, with equal attention for both positive and negative aspects of 
the applied method;
\item the development of a methodology based on this evaluation that is 
fit for use in corpus-based grammatical and linguistic studies;
\item an analysis and evaluation of the available tools in the field of 
papyrology for such studies, followed by suggestions for possible 
improvements to these tools;
\item the study of select grammatical and linguistic questions 
concerning the language of the Greek papyri using the methodology and 
toolset developed previously — essentially a practical test of the 
previous findings.
\end{enumerate}


%*******************************************************
% Acknowledgements
%*******************************************************


\begingroup
\let\clearpage\relax
\let\cleardoublepage\relax
\let\cleardoublepage\relax

\chapter{Acknowledgements}
\label{chap:acknowledgments}
\mtcaddchapter
I owe profound thanks to several people for the following
work. 

Firstly, I thank my supervisor, dr. Toon Van Hal, who from the get-go
demonstrated a very open-minded attitude in the face of unorthodox
subject matter for a classical philology thesis. He has demonstrated
exemplary patience with the stop-start rhythm of development of this
work. His corrections, suggestions and occasional nudges (usually
delivered electronically and with utmost tact) were an invaluable help
and encouragement. It is evident that without him, this work would not
have been possible.

Next, I am grateful to dr. Francis Maes, until recently a postgraduate
student at the Department of Computer Science at the KU Leuven Faculty
of Science, who during a conversation on this thesis directed me to a
set of state-of-the-art papers on machine learning for natural
language processing. His suggestions form the methodological bedrock
of this work in its current form.

I would also like to extend my thanks to my godmother Mar\'ia
Aguililla, who immediately accepted to proofread this work and did so
diligently, as well as Erwin De Koster for temporarily providing me
with computational power in the form of an iMac.

I offer my gratitude to my high school classics teacher, Guy van den
Heule, whose passion for Greco-Roman language and culture combined
with a lively and engaging teaching style pushed me to pursue the
study of the classical languages at a higher level.

I would also like to thank all my friends who offered me support and lended
me their proverbial ear in some way when the going got rough or showed
more than a cursory interest in this work, most of whom definitely
know who they are; they are too many to mention and I fear forgetting
but one.

I have not mentioned the most important persons of all: my parents,
who have been unflagging in their encouragement, interest and
support. They have enabled me to continue work on this thesis and to
pursue my study of computer science, which has granted me the
technical knowledge necessary to complete this work. To them I owe the
greatest thanks of all.

Finally, I invoke the memory of my good friend Lorena Lage Pi\~neiro,
who would have been mentioned in the previous paragraphs, if not for
the intervention of fate. Her sudden passing a week before this
writing deeply saddened me, and it is to her that I dedicate this
work.

\endgroup




\clearpage
%*******************************************************
% Contents
%*******************************************************
\phantomsection
\setcounter{tocdepth}{2}
\begingroup 
    \let\clearpage\relax
    \let\cleardoublepage\relax
    \let\cleardoublepage\relax

\dominitoc\tableofcontents
\let\clearpage\relax
\listoffigures
\listoftables
\endgroup
\markboth{\spacedlowsmallcaps{\contentsname}}{\spacedlowsmallcaps{\contentsname}} 

\begingroup 
    \let\clearpage\relax
    \let\cleardoublepage\relax
    \let\cleardoublepage\relax
\endgroup

\incrementmtc
\incrementmtc
\cleardoublepage

%******************************************************************
% main matter
%******************************************************************
\mainmatter
\pagenumbering{arabic}

%************************************************
\chapter{Introduction and preliminaries}
\label{chp:introduction}
%\minitoc\mtcskip
%************************************************

The study of the language of the papyri has in the past thirty years
seen little evolution until the recent appearance of Evans and Obbink's
\textit{The Language of the Papyri} \citep{lpapyri}, which has placed
the subject in the spotlight again. Twentieth-century scholarship on the
topic, though still useful for those interested in the study of the
papyri for historical purposes, is either antiquated, limited in scope
or incomplete (see \textit{infra}). Despite this, the papyri are
useful source material for the history and evolution of the Greek
language, as they contain not only official texts but private
documents as well, whose linguistic features and peculiarities have the
potential to foster new insights into the nature of colloquial Greek. 

\section{Thesis}
The following thesis intends to prove that it is possible to generate
basic linguistic annotation for a large digitalised corpus of papyri in
ancient Greek using readily available tools and techniques with minimal
technical overhead. Such a corpus could be  a boon to scholars interested
in the Greek of the papyri, as it would facilitate, for instance, the
creation of linguistically sound grammars and lexica.

\section{Preliminaries}

The following section will provide a background sketch, consisting of a short
overview of previous efforts and an elucidation of some key concepts.

\subsection{The language of the papyri}

The papyri began to be studied linguistically not by papyrologists and
historians, but rather by Bible scholars and grammarians interested in their
relevance in the development began to koin\^{e} Greek, particularly that of the
New Testament. G. N.  Hatzidakis, W. Cr\"onert, K. Dieterich, A. Deissmann, and
A.  Thumb pioneered the field in the late nineteenth and early twentieth
century, spurring a resurgence of scholarship on the topic; an excellent
overview of pre-1970s research may be found in \citet{mandilaras1973} and
\citeauthor{gignac1976} [\citeyear{gignac1976} and \citeyear{gignac1981}].

During this period, Mayser begon work on the earliest compendious grammar of
the papyri; it limits itself to the Ptolemaic era but explores it at length and
in great detail.  The work consists of a part on phonology and morphology, made
up of three slimmer volumes, and a part on syntax, encompassing three larger
volumes. Its composition seems to have been exhausting: it took Mayser
thirty-six years to finish volumes I.2 through II.3, with I.1 only completed in
1970 by Hans Schmoll, at which point the entire series was given a second
edition.

When casually browsing through some of its chapters (though casual is hardly
the word one would associate with the \textit{Grammatik}) it is remarkable to see
that Mayser brings an abundance of material to the table for each grammatical
observation he makes, however small it may be. For instance, the section on
diminutives essentially consists of pages upon pages of examples categorised by
their endings.

This is its great strength as a reference work - whenever one is faced with an
unusual grammatical phenomenon in any papyrus, consulting Mayser is bound to
clarify the matter; or rather, it was, for the work is now inevitably dated.
The volumes published during Mayser's lifetime only include papyri up to their
date of publication; only the first tome by Schmoll includes papyri up to 1968.
It is still a largely useful resource, but it is in urgent need of refreshment.

After Mayser set the standard for the Ptolemaic papyri, a grammar of the
post-Ptolemaic papyri was the new \textit{desideratum} in papyrology. The work
had been embarked on by Salonius, Ljungvik, Kapsomenos, and Palmer, only to be
interrupted or thwarted by circumstance or lack of resources.
\citet{salonius1927}, for instance, only managed to write an introduction on
the sources, though he offered valuable comments on the matter of deciding how
close to spoken language a piece of writing is. \citet{ljungvik1932} contains
select studies on some points of syntax.

It is in the 1930's that we see attempts to create a grammar of the papyri that
would be the equivalent of Mayser for the post-Ptolemaic period.
\citeauthor{kapsomenos1938} published a series of critical notes
[\citeyear{kapsomenos1938}, \citeyear{kapsomenos1957}] on the
subject; though he attempted at a work on the scale of the \textit{Grammatik},
he found the resources sorely lacking, as the existing editions of papyrus
texts could not form the basis for a systematic grammatical study. The other
was \citeauthor{palmer1934}, who had embarked on similar project and had
already set out a methodology [\citeyear{palmer1934}]; the war interrupted his
efforts, and he published what he had already completed, a treatise on the
suffixes in word formation [\citeyear{palmer1945}].

A new work of some magnitude presents itself two decades later with B. G.
Mandilaras' \textit{The verb in the Greek non-literary papyri}
[\citeyear{mandilaras1973}]. Though it does not aim to be a grammar of the
papyri, it does offer a thorough and satisfactory treatment of the verbal
system as manifest in the papyri.  Further efforts essentially do not appear
until the publication of Gignac's grammar. It is essentially treading in the
footsteps of Mayser, only with further methodological refinement and a more
limited, though still sufficiently exhaustive, array of examples. The author,
for reasons unknown to me, only managed to complete two of the three projected
volumes, on phonology and on morphology. The volume on syntax is thus absent, a
gap only partly filled by Mandilaras' \textit{The verb in the Greek
non-literary papyri}.

Finally, there is the aforementioned \textit{The Language of the Papyri}
\citep{lpapyri}, which does not aim to be a work on the same scale as the
aforementioned. It is a collection of articles on various topics, the whole of
which is meant to illuminate new avenues for future research. A particularly
relevant chapter for this thesis is the last one by Porter and O'Donnell
\citep{porter2010}, who set out to create a linguistic corpus for a selection
of papyri; their tagging approach, however, is manual, and their target corpus
limited. The authors also are the creators of \url{http://www.opentext.org/}, a
project aiming for the development of annotated Greek corpora and tools to
analyse them; sadly, no progress seems to have been made since 2005.

\subsection{Corpus linguistics}
A\footnote{The following section is based \emph{passim} on
\citet{okeeffe2010}.} corpus or text corpus is a large, structured collection
of texts designed for the statistical testing of linguistic hypotheses. The
core methodological concepts of this mode of analysis may be found in the
concordance, a tool first created by biblical scholars in the Middle Ages as an
aid in exegesis. Among literary scholars, the concordance also enjoyed use,
although to a lesser degree; the eighteenth century saw the creation of a
concordance to Shakespeare.

 The development of the concordance into the modern corpus was not primarily
 driven by the methods of biblical and literary scholars; rather, lexicography
 and pre-Chomskyan structural linguistics played a crucial role.

 Samuel Johnson created his famous comprehensive dictionary of English by means
 of a manually composed corpus consisting of countless slips of paper detailing
 contemporary usage. A similar method was used in the 1880s for the Oxford
 English Dictionary project - a staggering three million slips formed the basis
 from which the dictionary was compiled.

 1950s American structuralist linguistics was the other prong of progress; its
 heralding of linguistic data as a central given in the study of language
 supported by the ancient method of searching and indexing ensures its
 proponents may be called the forerunners of corpus linguistics.

Computer-generated concordances make their appearance in the late 1950s,
initially relying on the clunky tools of the day - punch cards. A notable
example is the \emph{Index Thomisticus}, a concordance
to the works of Thomas of Aquino created by the late Roberto Busa S.J. which
only saw completion after thirty years of hard work; the printed version spans
56 volumes and is a testament to the diligence and industry of its author. The
1970s brought strides forward in technology, with the creation of computerised
systems to replace catalogue indexing cards, a change that greatly benefited
bibliography and archivistics.

 It is only in the 1980s and 1990s that are marked the arrival of fully
 developed corpora in the modern sense of the word; for though the basic
 concepts of corpus linguistics were already widely used, they could not be
 applied on a large scale without the adequate tools. The rise of the desktop
 computer and the Internet as well as the seemingly ever-rising pace of
 technological development ensured the accessibility of digital tools.  The old
 tools - punch cards, mainframes, tape recorders and the like - were gladly
 cast aside in favour of the new data carriers.

 The perpetual increase of computing power equally demonstrated the limits of
 large-scale corpora; while lexicographical projects that had as their purpose
 to document the greatest number of possible usages could keep increasing the
 size of their corpora, the size of others went down as they whittled the data
 down to a specific set of uses of language.

 The possible applications of the techniques of corpus linguistics are diverse
 and numerous; for they allow for a radical enlargement in scope while
 remaining empirical, and remove arduous manual labour from the equation.
 Corpus linguistics can be an end to itself; it can, however, assert an
 important role in broader research.  \citet[7]{okeeffe2010} mention areas such
 language teaching and learning, discourse analysis, literary stylistics,
 forensic linguistics, pragmatics, speech technology, sociolinguistics and
 health communication, among others.

The term `corpus' has a slightly different usage in classical philology: they
designate a structured collection of texts, but that collection is not
primarily intended for the testing of linguistic hypotheses. Instead, we have,
for instance, the ancient corpus Tibullianum, or modern-day collection, for
  instance the Corpus Papyrorum Judaicarum, etc. We are primarily interested in
  the digital techniques used to create linguistic corpora; so let us first
  take a look at the progress of the digital classics.


\subsection{The digital classics} 

Classical philology, despite its status as one of the oldest and most
conservative scientific disciplines still in existence today, has in the last
fifty years found itself at the front lines of the digital humanities movement.
Incipient efforts in the fifties and sixties, mainly stylometric and lexical
studies and  the development of concordances, demonstrated the relevance of
informatics in the classics, an evolution that was at first met with some
skepticism, but later fully embraced.

The efforts began with the aforementioned Index Thomisticus, the first
computer-based corpus in a classical language; but the first true impetus was
the foundation of the Thesaurus Linguae Graecae project in 1972, a monumental
project with as its goal the stocking of all Greek texts from the Homeric epics
to the fall of Constantinople. Over the years, many functions have been added
to this ever more powerful tool; and even in the beginning stages of its
development, the TLG garnered praise.

The usefulness of the tool in its current form cannot be overstated: not only
does it contain a well-formatted and easily accessible gigantic collection of
text editions whose scope and dimensions exceed those of nearly any university
library; it also offers all of these texts in a format that allows for lexical,
morphological and proximity searches, as well as including a full version of
the Liddell \& Scott and Lewis \& Short dictionaries. The TLG has become a
staple of the digital classics.

Despite this, the TLG is becoming more and more dated as technology progresses.
While recent years have seen the rise of Unicode as the standard for encoding
ancient Greek, the TLG still uses beta code, a transliteration system designed
to only use the ASCII character set, and the texts are stored using an obsolete
text-streaming format from 1974, which divides the text in blocks of eight
kilobytes and marks the division between segment. 

A digitised version of the Liddell-Scott-Jones lexicon has been added to the
TLG's web interface, but the texts themselves have not undergone extensive
tagging, only lemmatisation.  Searching through the database can be done by
searching for specific forms of a lemma, or by searching for all forms of a
lemma, but this is essentially the limit of the search tool's power; it is not
possible to perform a query for all possible lemmata associated with a
particular form, i.e.\ we cannot find all forms which are, for example, an
active perfect indicative.

In the wake of the TLG, several notable projects have emerged: Brepols' Library
of Latin Texts is trying hard to be for Latin texts what the TLG is for Greek
texts; the Packard Humanities Institute has released CD's containing a
selection of classical Latin works. In more recent times, the Perseus Project
has enjoyed great popularity because of the attractive combination of an
excellent selection of classical texts with translations, good accessibility
and a set of interesting textual tools, the entire package carrying a very
interesting price tag for the average user — it is free to use, and for the
greatest part, open source as well.

The databases I have mentioned are quite general in scope; but within the
domain of classical philology, other specialised projects exist. Within the
field of papyrology, for instance, the digital revolution has taken a firm
foothold. Starting with several separate databases, the field has experienced a
tendency towards convergence and integration of the available resources, as
exemplarised by the papyri.info website, maintained by Columbia University,
that integrates the main papyrological databases into a single database.

A great feature of this database is the shell in which all data is wrapped;
they are compliant with the EpiDoc standard, a subset of XML based on the TEI
standard and developed specifically for epigraphical and papyrological texts.
One may access the database’s resources through the Papyrological Navigator and
suggest corrections and readings through the Papyrological Editor. What’s more,
all data is freely accessible under the Creative Commons License,
crowd-sourced, regularly updated, and can be downloaded for easier searching
and tweaking.

In other words, papyri.info has brought the open-source mentality from the
computer world into the classics. For our purposes, this open setup is
desirable, as the database is not fit for them as it is, but can with some
effort be molded into a useful tool.

\subsection{Natural language processing} 

Natural language processing (henceforth NLP) is a subdiscipline in computer
science concerned with the interaction between natural human language and
computers. Its history well and truly starts in the fifties, with a basic
concept which has played a great role in natural language processing, and
computer science in general, the Turing test. This test, put forth by Alan
Turing in his seminal paper \textit{Computing Machinery and Intelligence}
\citep{turing1950}, evaluates whether a machine is intelligent or not by
placing a human in conversation with another human and a machine; if the first
human cannot tell the other human and the machine apart, the machine passes the test.

Machine translation systems entered development, though progress soon stalled
because of technical limitations and because of methodological obstacles: such
systems were dependent on complex rulesets written by programmers that allowed
for very little flexibility. Because of the slow return on investments made,
funding for artificial intelligence in general and machine translation
specifically was drastically reduced throughout the late sixties and the seventies.

A resurgence followed: thanks to advances in computational power and the
decline of Chomskyan linguistics, which had been the dominant theoretical
vantage point in the preceding thirty years, the eighties were marked by the
introduction of statistical machine translation, which is fundamentally based
on the tenets of corpus linguistics. Modern natural language processing is
therefore situated on the crossroads between various fields: artificial
intelligence, computer science, statistics, and corpus and computational
linguistics. It looks to be an exciting field for the coming years as its
techniques are under constant improvement and ever more present in our daily lives.

Most NLP software is designed explicitly with living languages in mind;
English, being a world language and the international \textit{lingua franca},
has enjoyed most of the attention, but other major languages have enjoyed some
attention, too. Ancient languages, however, are neglected, presumably due to
their often high complexity and the extensive study and analysis to which they
have been submitted by skilled scholars. Yet most texts have not been
integrated in annotated corpora; and though databases such as the Perseus
project contain large swathes of morphologically and sometimes syntactically
annotated text, the process has been driven largely by manual labour; to give
an exhaustive list is not appropriate here, but another such example which is
relevant is the PROIEL project, which is also a treebank, i.e.\ a database of
syntactically annotated sentences. It contains data for Herodotus and the New Testament.

On the other hand, there are also corpora which have been tagged using NLP
techniques, whose relevancy for this thesis is high and that I have thus
described in the next section on methodology.

\section{Methodology}

In this thesis, we have been largely inspired by two articles by H. Dik and R.
Whaling, [\citeauthor{dik2008}, \citeyear{dik2008} and \citeyear{dik2009}], in
which they document their method for semi-automatically tagging the Perseus
project's texts under their own framework, PhiloLogic. They start with a
database of analysed forms and a series of tagged texts which they use as
initial data to train a decision tree tagger, TreeTagger, developed by Helmut
Schmid at the University of Stuttgart, a tool which despite being developed in
1995 has aged well as far as performance is concerned. They achieved remarkable
accuracy: with refinements to the training data they achieved 96.2\% accuracy
during tests on the original training data and 91\% accuracy on new data, a
result which compares quite favorably when compared to TreeTagger's 97\%
accuracy when used on German newspaper articles considering the high complexity
of ancient Greek and the variety of styles of ancient Greek literature.

It occurred to me that this might be a great method for processing the
papyri.info database with a relatively small effort for a high payoff; using
data from the Perseus and PROIEL projects, it could be possible to train
TreeTagger for both morphology and syntax, apply the resulting parameters to
the corpus and thus for the greatest part obviate the need for manual tagging.
Given the extent of the corpus (about 50,000 texts containg almost 4,500,000
words), achieving even 85\% accuracy would reduce the amount of untagged words
to 675,000, many of which I would expect to be proper names or morphologically
`erroneous' forms as are often found in the papyri, data which could itself be
analysed with regular expressions and then used to improve the training
data.\footnote{As I set out to verify the originality of my thesis, I found
  that this statistical approach has been used before for textual criticism!
  \textit{Vide} \citet{mimno2009}, an abstract of which may be found at
  \url{http://people.cs.umass.edu/~wallach/publications/mimno09computational.txt}.}

%************************************************
\chapter{The corpus}
\label{chp:design}
\minitoc\mtcskip
%************************************************
\section{Goals}

Our objectives for the modified corpus is threefold. Firstly, we want to
morphologically annotate it with reasonable accuracy; secondly, we also want to
add syntactical annotation; thirdly, we want to ensure the corpus is compatible
with a broad range of tools.

The former two goals are the central enterprise of this thesis. Morphological
annotation will pose little theoretical problems, as the morphological features
of Greek are not subject to discussion; but for syntactical annotation, a
variety of theoretical frameworks are in existence. The fact of the matter is
that I have chosen to take the route which is technically easiest: to use the
tenets of dependency grammar. The reason Leuven : K.U.Leuven. Faculteit Letteren, 2011 is twofold: on one hand, our training
data, the New Testament annotated by PROIEL, itself is structured using a
dependency model; on the other hand, dependency trees, in contrast to
constituent analysis, assigns each word or morph one node, which is easier to
`digest' for a computer program.

Our third goal is to format the corpus in a way that will ensure broad
compatibility. This is desirable for several reasons. A first is that a broadly
used format offers many options for data treatment; for instance,
comma-separated text files can be read by many programs, from the simplest text
editor to the most complex software packages, while XML can very easily be
integrated in a web interface or parsed and transformed by a variety of
programs. A second is that corpus is to be published on GitHub for
collaborative editing, or possibly pulled into the main papyri.info repository.
This depends on our first point: the idea of GitHub is that all projects are
open source and can be contributed to by others, which is facilitated by using a
format with wide adoption.

\section{Design}
\subsection{Initial inquiries}

Before arriving at the approach described above, I attempted several other
possible methods for automatic analysis. A first idea was to integrate the XML
files in a PhiloLogic setup.  PhiloLogic is a tool developed by the ARTFL
project and the Digital Library Development center at the University of Chicago
and released under the GNU General Public License. Its original purpose was to
serve as the full-text search, retrieval and analysis tool for large databases
of French literature.

PhiloLogic has support for TEI XML and boasts an impressive
array of features: it can search through text and deliver KWIC results
filtered by frequency and metadata, as well as collocation tables and word
order information.  Furthermore, the development web page cites support for
bibliographic backends in MySQL databases, out of the box operation,
interoperability across Unix-based systems, etc.

The tool seemed promising; especially the built-in support for XML and
Unicode drew my attention, along with the built-in features for linguistic
analysis. It essentially would have made my work easier and allowed me to
invest more time in investigating a few linguistic topics. In the end, however,
the software did not fit my needs in a few respects.

Firstly, PhiloLogic requires Apache, Perl and several CPAN modules, as well as
gawk, gdbm, and agrep.  Not a huge amount of dependencies, but still less than
using Python and the NLTK, which does not require a server or SQL database to
run; also a big problem is the out-of the box incompatibility of the required
CPAn  modules with 64-bit systems. Setting PhiloLogic up was not as
easy as advertised; I had to resort to a virtual machine running 32-bit Debian
Linux until I was able to discover a method that enabled compatibility with
Mac OS X 10.7, my OS of choice.

Secondly, XML support includes standard TEI  but has issues with EpiDoc; notably,
in my experience, headers were treated as text and EpiDoc's complicated tagset
was incorrectly rendered; for instance, original readings and corrections were
concatenated in the text browser, an undesirable situation if we wish to have
any hand in our choice of corpus, and a possible source of statistical errors.
Developing brand new XSLT stylesheets to convert EpiDoc to a simpler form of
TEI markup or to raw text could perhaps have been a feasible solution; but
PhiloLogic seemed to behave erratically in general and a first attempt at
feeding the system raw text (advertised as one of its capabilities) turned out
a spathe of errors.

Thirdly, PhiloLogic's feature set is extensive, but I soon came to realize that it
does not offer much more than papyri.info already does; it is essentially a
robust concordancing application through a web interface. I wished not only to
have a lemmatised corpus searchable by word proximity, but also by word
relations, which would shed more light on the syntax of the Greek of the
papyri, as this is arguably the largest gap in scholarship due to the age of
Mayser's \emph{Grammatik} and the incomplete state of Gignac's \emph{Grammar}.
Furthermore, the tool as available to the public does not integrate the
morphological database of Greek developed for Perseus under PhiloLogic --- a
regrettable point, really, since it was exactly this capacity I wished to
exploit in the first place. Unicode search for Greek also seems to be in need
of improvement; the search form still requires transcription. Modifying the
tool would have been unfeasible given the size of the source code.

Thus I discarded PhiloLogic from my options and set out to find or develop a
simpler system.\footnote{I also considered setting up a mirror of papyri.info
on a personal server and modifying the search functionality; this turns out not
to be a task for the weak-hearted. The required setup causes substantial
overhead for our purposes, too much to be a reasonable solution.} After the
technical struggle with PhiloLogic, I decided to focus my efforts on
concordancers that could replicate its functions. A quick consultation of
Stanford's list of resources \citep{stanfordnlpres} for statistical natural
language processing and corpus-based computational linguistics directed me to
WordSmith Tools, which is regrettably Windows-only and a commercial, closed
source program to boot.  I set out on a search for open-source alternatives and
found a diversity of programs, most of them without the same functionality as
WordSmith, and invariably clunky or antiquated --- making WordSmith the only
option, but even that did not seem viable for my ends.

I continued my search and stumbled across a very interesting piece of software
that is enjoying a good amount of popularity as a didactic tool for
computational corpus linguistics; the Python NLTK, short for Natural
Language Tool Kit, is not a stand-alone application, but rather a set of ``open
source Python modules for research and development in natural language
processing and text analytics bundled with data and documentation''
\citep{nltkhome}. These modules implement diverse functionalities useful for
natural language processing and allow their easy integration into Python
programs. The Python NLTK  seemed to offer a more than solid amount of
features: corpus readers, tokenizers, stemmers, taggers, chunkers, parsers,
classifiers, clusterers, tools for semantic interpretation and metrics were all
integrated from the get-go and are easily combined and extended. The
NLTK was developed for English, but its extensibility meant it could
be applied to ancient Greek with relative ease. A further advantage is the fact
that Python is very well-suited to text manipulation and parsing XML,
which is an advantage when working with the papyri.info corpus; it is also
cross-platform: Python interpreters are available for all major platforms.

Now that I had a solid tool for interpreting corpora, the only issue remaining
was to construct a linguistically annotated corpus out of the papyri.info
corpus. The extent of the corpus makes extensive manual tagging by a single
person unfeasible; thus I set out to look for an automated tagger which could
be used, immediately or with some effort,  on ancient Greek. A variety of
taggers seem to be available, but I settled on TreeTagger, developed by Helmut
Schmid, for several reasons:

\section{Technical} % (fold)
\label{sec:corpus-technical}

\subsection{Requirements} % (fold)
\label{sub:requirements}

To create the annotated corpus as I have, it is necessary to have the following installed (older versions may work but have not been tested):

\begin{itemize}
  \item a UNIX-like operating system, i.e. Mac OS X, any variety of Linux or BSD (using Windows should also be possible, since all tools used are portable and cross-platform, but the file location hierarchies in some scripts will not work);
  \item Python 2.7.3, found at \url{http://www.python.org/} (do not use Python 3.x, the NLTK is not compatible with it);
  \item Perl 5.x \url{http://www.perl.org/};
  \item Saxon-HE 9.x, found at \url{http://saxon.sourceforge.net/};
  \item Helmut Schmid's TreeTagger, found at \url{http://www.ims.uni-stuttgart.de/projekte/corplex/TreeTagger/};
  \item the Python Natural Language Toolkit (NLTK), found at \url{http://nltk.org/};
  \item SQLite 3, found at \url{http://www.sqlite.org/}.
\end{itemize}

All the above, excluding Mac OS X, are freely available.

% subsection Requirements (end)

\subsection{Extracting training data} % (fold)
\label{sub:extract_training_data}

% subsection Extracting training data (end)

\subsection{Principles of statistical natural language processing} % (fold)
\label{sub:principles-nlp}

This subsection is intended as an overview of some key concepts on which
natural language processing is based; it is based mainly on 
\citet{manning1999}.

\subsubsection{Formal languages} % (fold)
\label{sub:formallang}

% subsection Formal languages (end)

\subsubsection{Hidden Markov models} % (fold)
\label{ssub:markovchains}

% subsubsection Markov chains (end)

\subsubsection{Zipf's law} % (fold)
\label{ssub:zipflaw}
% subsubsection Zipf's law (end)

% subsection Principles of statistical natural language processing (end)

\subsection{Accuracy testing} % (fold)
\label{sub:accuracy-testing}

% subsection Accuracy testing (end)

% section Technical (end)
%************************************************
\chapter{The tool}
\label{chp:design}
\minitoc\mtcskip
%************************************************
\section{Goals} % (fold)
\label{sec:tool-goals}

% section Goals (end)

\begin{enumerate} 
  \item The tool must be able to read EpiDoc XML and a subset of
      additional tags added for linguistic annotation. It should be,
      essentially, a text-based replica of the Papyrological Navigator as found
      at \texttt{papyri.info}. 
    \item search 
    \item statistics
    \item syntactic trees
\end{enumerate}

\section{Design} % (fold) 
\label{sec:tool-design} 
The tool works from the command line; while it has the obvious disadvantage of
being a type interface most people cannot work with right off the bat, command
line programs generally have an edge over programs with a graphical user
interface: they are smaller, faster, and more efficient, and in Unix
environments can easily be integrated in a larger workflow.
% section Design (end)

\section{Technical} % (fold)
\label{sec:tool-technical}
The tool is written in Python 3. 
% section Technical (end)

\section{Interface} % (fold)
\label{sec:interface}

The basic command for launching the script is \texttt{python tjufy} from within
a terminal or simply \texttt{tjufy} from within the Python interpreter (from
here on I will simply write \texttt{tjufy}). Running it without options like this
will display a help dialog; the same goes for running \texttt{python tjufy -h}
or \texttt{python tjufy --help}.


% section User interface (end)

\chapter{Applications} % (fold)
\label{cha:applications}

\section{Collaborative editing} % (fold)
\label{sec:collaborativeediting}

% section Collaborative editing (end)

\section{Corpus-based grammars} % (fold)
\label{sec:corpusbasedgrammars}

% section Corpus-based grammars (end)

\section{Textual criticism} % (fold)
\label{sec:textualcriticism}

% section Textual criticism (end)

% chapter Applications (end)

\chapter{Conclusion} % (fold)
\label{cha:conclusion}

% chapter Conclusion (end)

% *****************************************************************
% back matter
%******************************************************************
\backmatter
\clearpage
%*******************************************************
% Bibliography
%*******************************************************
\nocite{*}
\printbibliography


%\clearpage
%% %************************************************
% \appendix
% \chapter{Appendix}
% \label{chap:appendix}
% %************************************************
\appendix
\pretocmd{\chapter}{%
  \clearpage
  \pagenumbering{arabic}%
  \renewcommand*{\thepage}{\thechapter-\arabic{page}}%
}{}{}

\chapter{Concepts and techniques}
\label{chp:conceptstechniques}
\section{Mathematics}
\label{sec:mathematics}


\subsection{Set theory}
\label{sec:settheory}

Consider an object $o$ and a set $A$. We write 'object $o$ is an
element of set $A$' as $o \in A$. Sets themselves are also objects and
can belong to other sets. Consider two sets, $A$ and $B$. We write
that $A$ is a subset $B$ as $A \subseteq B$; this implies that $B$
contains all elements in $A$ and does not exclude the possibility of
$A$ = $B$; if $A$ is a proper subset of $B$, i.e. all elements of $A$
are in $B$ but not all elements of $B$ are in $A$, we write $A \subset
B$. Mirroring these symbols from right to left gives us the symbols
for supersets and strict supersets, respectively. The empty set, which
contains no elements, is written as $\varnothing$.

There exist different binary operators on sets (i.e. operators on two
sets) which return another set. The most common operators are:
\begin{itemize}
\item the union of sets, written as $A \cap B$, denotes a set which
  contains all elements which are in $A$ and $B$;
\item the intersection of sets, written as $A \cup B$, denotes a set
  which contains all elements which are in both $A$ and $B$;
\item the difference of sets, written as $A \setminus B$, denotes a
  set which contains every element from $A$ excluding those which are
  also in $B$.
\item the Cartesian product of sets, written as $A \times B$, is the
  set containing all ordered pairs of elements from $A$ and $B$.
\end{itemize}

\subsection{Probability}
\label{sec:probability}

A \textbf{probability} is a measure for the likelihood of an event for an
experiment, which we intuitively understand to be an action whose
outcome we want to observe. Such events are then elements of a set
containing all possible outcomes of an experiment; an event can be a
point or subset of that set. We call this set the \textbf{sample space},
denoted S. We denote the probability of an event E as P(E).

Axiomatically, we can define probability as follows:

\begin{enumerate}
\item For every event $E$, $P(E) \geq 0$; no event can have a negative probability.
\item $P(S) = 1$; that is to say, every experiment has an event.
\item For any sequence of disjoint events $A_i$ (that is to say,
  there is no overlap between events), the probability of any one of
  these events occurring is the sum of their respective probabilities.
\end{enumerate}

A few other important properties of probabilities and events are the following:

\begin{enumerate}
\item The complement of an event $A$, which is the union of all
  elements of $S$ which are not an element of $A$, is denoted $A^C$;
  $P(A^C)$ is the probability of this event occurring and is equal to $1
  - P(A)$.
\item For any event E, $0 \leq P(E) \leq 1$.
\item Given any two events $A$ and $B$, $A \subset B \rightarrow P(A) \leq P(B)$.
\end{enumerate}

Given probabilities of a number of events, we can establish
relationships between these probabilities and compute related
probabilities using certain rules. A classic rule is the
\textbf{multiplication rule}, which states that if we perform an
experiment in $k$ parts and the $i^{th}$ part of the experiment has
$n_i$ possible outcomes, and the outcomes of prior parts of the
experiment do not affect latter ones, the probability of any specific
sequence of partial outcome will be the product of all outcome counts
$n_i$ with $i$ ranging from $1$ to $k$.

Set theory is important when we want to know the probability of an
event E which can be constructed from a set of sets $A_i$ using set
operators. The third axiom of probability has already given us the
solution for disjoint events; events may also overlap, and in this
case, we need a more sophisticated formula. For the union of any $n$
events $A_i$, the following holds:
\begin{equation}
  P(\bigcup_{i=1}^n A_i) = \sum\limits_{i=1}^n P(A_i)
\end{equation}

For the intersection of these $n$ events $A_i$, it can be proven that:
\begin{equation}
  P(\bigcap_{i=1}^n A_i) = \lim_{n \to \infty} P(A_i)
\end{equation}
Knowing both these rules is important when considering
\textbf{conditional probability}; for two events $A$ and $B$, suppose
we know $B$ has occurred and we want to know what the probability of
$A$ occurring is given this information. We call this the conditional
probability of $A$ given $B$ and write it $P(A|B)$. If $P(B) > 0$,
then we define it as:

\begin{equation}
  P(A|B) = \frac{P(A \cap B)}{P(B)}
\end{equation}

Using this formula, we can directly derive \textbf{Bayes' theorem} as follows:

\begin{equation}
  \begin{aligned}
    P(A|B) &= \frac{P(A \cap B)}{P(B)} \\
    P(A|B) P(B) &= P(A \cap B)\\
    P(A|B) P(B) &= P(B \cap A)\\
    P(A|B) P(B) &= P(B|A) P(A)\\
    P(B|A) &= \frac{P(A|B) P(B)}{P(A)}
  \end{aligned}
\end{equation}

\subsection{Calculus and linear algebra}
\label{sec:calculus}

\subsubsection{Derivates and extrema}
The derivative of a function is a measure of its rate of change. For a
function in one variable, we can visualise the derivative as a
function which for each point gives the slope of the tangent line to
the original function. Derivatives can be used to find the extrema of
a function, that is to say, the 'low' and 'high' points, with the help
of the fact that the tangent line at these extrema has slope zero. By
'following' the direction indicated by the derivative, we can find
such extrema.

Derivatives are defined for functions with any number of variables,
and so are extrema. The derivative of a multivariable function is
called a gradient; this is a vector with each component containing the
partial derivative of the function with respect to a single of its
variables. We can find extrema for multivariable function by applying
the same method as we do for functions in one variable, by 'following'
the gradient. This is called gradient descent or ascent, depending on
whether we want to find a minimum or a maximum for a function.

\subsection{Statistics}
\label{sec:statistics}
    \subsubsection{Regression}

    Regression is a classic technique from statistics, often visualised as
    'fitting a line to a set of points'. Classification is a related
    technique which uses regression to classify new data points. This is
    essentially what any probabilistic model for natural language
    processing does, in one form or another. What follows is an overview
    of various types of regression and corresponding methods for
    classification.

    We start with \textbf{univariate linear regression}. Given a set
    of $n$ points in the plane, we want to find a hypothesis that best
    corresponds to the location of these points, and we want this
    hypothesis to be a linear function. This function is then of the form:
    \begin{equation}
      h_w(x) = w_1x + w_0
    \end{equation}
    Unless all points are collinear, it is of course impossible to
    find a function of this form that gives a correct mapping for each
    point. The best we can do is find the values of $w_o$ and $w_1$ for
    which the empirical loss on the mappings is minimal. The traditional
    way of doing this is to define a function that computes the squares of
    the errors and sums it over all data points; this is called an
    \textbf{$L_2$ loss function}. We now want to find the values of $w_0$
    and $w_1$ for which this function attains a minimum. We can find these
    minimal points by solving for the roots of the partial derivative
    functions of this loss function with respect to $w_0$ and $w_1$.  This
    problem is mathematically relatively simple and has a unique
    solution. This solution is valid for all loss functions of this type.

    Problems arise when we are trying to create a nonlinear model. In
    this case, the minimum loss equations frequently do not have a unique
    solution. We can of course still model the problem algebraically, and
    the goal is the same: finding the roots of the partial derivative
    function. Now, however, we need to use a more sophisticated method:
    \textbf{gradient descent}. We can visualise this technique as
    'descending a hill'; the 'hill' is the graphical representation of the
    root of the system of partial derivatives, and by 'descending' this
    hill, i.e. by iteratively picking values which bring us closer to the
    bottom part of the valley next to the hill, which corresponds to the
    minimal point of the function, eventually convergence will be reached
    on the minimum and we will found the correct weights for our
    function. The difference by which we change the value at each
    iteration is called the \textbf{step} or \textbf{learning rate} and
    determines how fast we will converge; it may be either a fixed
    constant or a mutable value which can increase or decay according to
    the current state of our descent.

    \textbf{Multivariate linear regression} poses a similar problem;
    only this time the function is not dependent on a single variable, but
    on two or more. Such a function is a bit more complex, but we can
    find a solution to the regression problem using analogous
    techniques. Suppose the function has $n$ variables. Each example $x_j$
    must be a vector with $n$ values. At this point, we are looking at a
    function of the following form:
    \begin{equation}
      h_w(x_j) = w_0 + w_1x_{j,1} + w_2x_{j,2} + ... + w_nx_{j,n} = w_0 + \sum\limits_{i} w_ix_{j,i}
    \end{equation}
    We want to simplify this to make algebraic manipulations
    easier. We therefore prepend an extra component $x_{j,0} = 1$ to the
    vector $x_j$; now using vector notation we can simplify the
    previous equation to:
    \begin{equation}
      h_w(x_j) = \sum\limits_{i} w_ix_{j,i} = w \cdot x_j
    \end{equation}
    What we are now looking for is a vector $w$ containing the weights
of our function which minimises the empirical loss, as in univariate
linear regression. We can equivalently use gradient descent; only now,
of course, the computational cost of that technique will be higher. A
common problem can now appear: \textbf{overfitting}, that is, giving
an irrelevant dimension of the vector $w$ too much weight due to
chance errors in the computation. This can be compensated by taking
into account the complexity of the hypothesis; a statistical
equivalent to Ockham's razor, if you will.

\subsubsection{Classification} 

We can define an analogous process for classification; only now
the function must not fit to the data itself but must create a
\textbf{decision boundary} between data points. If there exists a
linear function which satisfies this property for a given data set
, we call the bounding line or surface generated by this function a
\textbf{linear separator}, and the data set \textbf{linearly
  separable}. The hypothesis function is now of the form:
\begin{equation}
  h_w(x) = 1$ if $ w \cdot x \geq 0$, $0$ otherwise.$
\end{equation}
We can view this as a function $threshold(w \cdot x)$ which is
equal to 1 only if $w \cdot x = 0$ Note that while the separating
function is linear, the hypothesis function is not, and in fact has
the distinctly unappealing property of not being differentiable. We
can therefore not apply the technique of gradient descent
here. Furthermore, this type of function has exactly two outputs: 1 or
0. For our purposes, we need subtler methods of classification.  This
type of hypothesis function is therefore not fit for our purposes, but
it does give a good idea of what classification is.

The best option is replacing the hard threshold function with the
sigmoid or logistic function, which offers a good approximation and is
differentiable at every point. This function is of the form:
\begin{equation}
  g(x) = \frac{1}{1 + e^{-x}}
\end{equation}
Such that our new hypothesis function is:
\begin{equation}
  h_w(x) = g(w \cdot x) = \frac{1}{1 + e^{- w \cdot x}}
\end{equation}
If we use this function, we are performing \textbf{linear
  classification with logistic regression}.

\subsubsection{Logistic regression and the chain rule}
\subsubsection{Clustering}

\section{Natural language processing}
\label{sec:nlp}

\subsection{N-grams}
\label{sec:ngrams}

\subsection{Hidden Markov Models}
\label{sec:hmm-maxent}

\subsection{Viterbi decoding}
\label{sec:viterbi}

\section{Artificial intelligence and machine learning}
\label{sec:aiml}

\subsection{What is machine learning?}
\label{sec:statistics}
\begin{itemize}
\item Supervised learning
\item Unsupervised learning
\end{itemize}

\subsection{Neural networks}
\label{sec:neuralnetworks}

An artificial neural network is a massively parallel distributed
processor made up of simple processing units, which has a natural
propensity for storing experiential knowledge and making it available
for use.

For the design of an architecture which allows us to solve the
problems above, we have taken our cues largely from recent work in
machine learning as applied to natural language processing. In
particular, we follow the approach set out in Weston \& Collobert 2008
and expanded in Weston et al. 2011, that is, the use of deep neural
networks for joint training on our chosen corpus. This section is
dedicated to a more expository overview of that architecture for the
mathematical layman.

An artificial neural network is a massively parallel processing system
constructed from simple interconnected processing units (called
neurons) which has the ability to learn by experience and store this
knowledge for later use. The term 'neural network' is due to the
resemblance of this architecture to the most powerful biological
processor known to exist, the human brain, which has a way of
functioning which is broadly analogous to this process. 

Artificial neural networks find their origin in a mathematical model
dating from before the first wave of artificial intelligence in 1956,
the McCulloch-Pitts Threshold Logical Unit, known also as the
McCulloch-Pitts neuron. Warren McCulloch was a psychiatrist and
neuroanatomist; Walter Pitts a mathematical prodigy. Both met at the
University of Chicago, where a neural modeling community led by the
mathematical physicist N. Rashevsky had been active in the years
preceding the publication in 1943 of the seminal paper A logical
calculus of the ideas immanent in nervous activity. 

Formally, we can define a neuron as a triplet (v,g,w) where:

\begin{itemize}
\item v is an input function which takes a number of inputs and computes their sum;
\item g is an activation function which is applied to the output of the input function and determines if the neuron 'fires';
\item w is an output function which receives its input from the activation
  function and distributes it over a number of outputs.
\end{itemize}

Given its structure, we can also see a neuron as a composite function
$F = w \circ g \circ v$. The combination of several of these units
using directed links forms a neural network. A link connecting a node
$i$ to a node $j$ transfers the output $a_i$ of node $i$ to node $j$
scaled by a numeric weight $w_{i,j}$ associated with that specific
link. This is the general model of a neural network; countless
variations on this theme have been developed for different purposes,
mainly by modifying the activation function and the interconnection of
neurons.

The activation function $g$ typically will be a hard threshold
function (an example of this is the original McCulloch-Pitts neuron),
which makes the neuron a perceptron, or a logistic (also known as a
sigmoid) function, in which case we term the neuron a sigmoid
perceptron. Both these functions are nonlinear; since each neuron
itself represents a composition of functions, the neuron itself is a
non-linear function; and since the entire network can also be seen as
a composite function (since it takes an input and gives an output) the
network can be viewed as a nonlinear function. Additionally, choosing
to use a logistic function as an activation function offers
mathematical possibilities, since it is differentiable. This offers
similar possibilities as the use of the logistic function for
regression (cf. \textit{supra}).

The links between nodes can be configured in different ways, which
each afford distinct advantages and disadvantages. Broadly, we can
distinguish two models. The simplest model is the \textbf{feed-forward
  network}, which can be represented as an acyclic directed graph. The
propagation of an input through this kind of network can be seen as a
stream, with posterior (downstream) nodes accepting outputs from prior
(upstream) nodes. This type of network is the most widely-used and is
used in the architecture. A more complex type is the \textbf{recurrent
  network}, which feeds its output back to its input and thus contains a
directed cycle; this type of network has interesting applications (for
example in handwriting recognition), as they resemble the neural
architecture of the brain more closely than feedforward networks do.

Feed-forward networks are often organised (to continue the stream
analogy) in a kind of waterfall structure using layers. The input is
the initial stream, the output is the final stream; in between, we may
place hidden layers, which are composed of neurons which take inputs
and outputs as any neuron does, but whose output is then immediately
transferred to a different neuron. Throughout the network, we can
equip the neurons in each layer with distinct activation functions and
link weights and in this way mold the learning process of the network
to our purpose.

Single-layer networks contain no hidden layers; the input is directly
connected to the output. Therefore, the output is a linear combination
of linear functions. This is undesirable in many cases. The main
problem, demonstrated early on in the development of neural network
theory, is the fact that such a network is unable to learn functions
that are not linearly separable; one such function is the XOR
function, which is a very simple logical operator. Despite this, such
neural networks are useful for many tasks, as they offer an efficient
way of performing logistic regression and linear classification.

Our interest lies in multi-layer networks, however. Multi-layer
networks contain one or more layer between the input and output layer,
which are called hidden layers. By cascading the input through all
these layers, we are in fact modeling a nonlinear function which
consists of nested nonlinear soft threshold functions as used in
logistic regression. The network can now be used to perform \textbf{nonlinear
  regression}. Different algorithms exist which can be used to train the
network; the most important one is the \textbf{backpropagation algorithm},
which is the equivalent of the loss reduction techniques used in
linear regression.

Suppose that a neural network models a vector-valued hypothesis
function $h_w$ which we want to fit to an example output vector
$y$. We can create a $L_2$ loss function E by taking the error on this
vector and squaring it. This function can be quite complex, but by
taking partial derivatives of this function, we can consider the
empirical loss on each output separately, like so:

\begin{equation}
  \begin{aligned} 
    \frac{\partial}{\partial w} E(w) &= \frac{\partial}{\partial w} \lvert y - h_w(x) \rvert^2 \\
    &= \frac{\partial}{\partial w} \sum (y_k - a_k(x))^2 \\
    &=  \sum \frac{\partial}{\partial w} (y_k - a_k(x))^2
  \end{aligned} 
\end{equation}

If the output function is to have $m$ outputs, instead of handling one
large problem, we can subdivide the problem into $m$ smaller
problems. This approach works if the network has no hidden layers, but
due to the fact that nothing is really known about the hidden layers
if we only look at the output layer, a new approach is necessary. This
is called backpropagation, a shorthand term for backward error
propagation.

Backpropagation relies on knowing all the network weights; from these,
given an error at the output layer, we can determine what part of the
error is due to which neuron in the hidden layer. We use this
knowledge to establish update rates for neurons in the preceding
layer. For each successive layer, working backwards from the output
layer, we can apply the same process and eventually make good
parameter adjustments.


\subsection{Deep learning}
\label{sec:deeplearning}

A single-layer neural network is a compact structure that can perform
complex tasks efficiently; and even with the hardware which was in use
a decade ago and before, training a neural network was a feasible
task. Training a deep neural network, which contains several hidden
layers (hence the term \textit{deep}), is another matter; doing this
requires the backpropagation algorithm, which until the mid-2000's was
simply too slow for use on the hardware of the day.

At this point, G. Hinton, who had been one of the first proponents of
the use of deep networks in the 1980's as a professor at Carnegie
Mellon University, blew new life into the field he helped create with
his paper ``Learning multiple layers of representation''
\citep{hinton2007learning}. He demonstrated that it was possible to
perform very complex learning tasks relatively efficiently and to
great effect using this type of network.

The way this is done is by stepping down from the traditional approach
where neural networks are given a certain amount of classified
examples and training them to classify observations based on these;
rather, the goal is now to train \textbf{generative networks}, i.e.,
networks that can randomly generate possible observations. Examples
given to the network do not need to be classified in advance; instead,
given an observation, the network's parameters are adjusted to
maximise the likelihood of data of its kind being
generated. 

A classic task for this type of network is handwritten digit
recognition; a network is trained by feeding it a large amount of
examples of handwritten digits. For each example, the parameters are
adjusted; after a sufficiently large amount of examples, the network
is capable of generating handwritten digits itself, in a vast number
of variations.


% \section{Principles of statistical natural language processing} 
% \label{sub:principles-nlp}
% \subsection{Context-free grammar and pushdown automata} % (fold)
% \label{sub:formallang}

% The basis of all statistical natural language processing is the
% conceptualisation of natural language within the Chomsky hierarchy as a formal
% language generated by a \textbf{stochastic context-free grammar}, that is to say, a
% context-free grammar whose production rules are augmented by a probability. In
% computer science, such a language is termed a context-free language and defined
% as any language that can be recognized by a \textbf{nondeterministic pushdown
% automaton}. A nondeterministic pushdown automaton is understood to be a
% nondeterministic finite state automaton with access to an infinite stack, a
% finite state automaton being an automaton which is in a state that can
% transition into another state when triggered by input. Put more simply, a
% nondeterministic pushdown automaton possesses the following:

% \begin{itemize} 

% \item a \textbf{stack}, that is, a string of symbols which functions as the
%   automaton's memory. The automaton only has access to the leftmost symbol in
%   this string, a symbol which is dubbed the top of the stack; the top of the
%   stack can be used to determine which transition the automaton will make
%   next;

% \item an \textbf{input tape} whose symbols are scanned one by one just like a
%   finite-state automaton; 

% \item a \textbf{finite state control unit}, which controls the state of the
%   automaton. Normally this is done in response to input, but it is possible
%   in some types of automata to transition into a new state without any
%   input, as is the case for pushdown automata; this type of transition is
%   termed $\epsilon$-transition or $\lambda$-transition.

% \end{itemize}

% A simple example of finite state automaton is a coin-operated turnstile: it has
% two possible states, locked and unlocked, and two possible inputs, the
% insertion of a coin or a push. Depending on the state the machine is currently
% in, each of these inputs will trigger either a change or leave the machine in
% its current state. For instance, if the turnstile is unlocked, a push will make
% it turn and lock again; inserting a coin will do nothing.  Conversely, when it
% is locked, a coin will unlock it, while a push will still leave it locked. This
% type of automaton is deterministic, that is, there is only one next possible
% state to transition to. A nondeterministic finite state automaton would for
% example be a vending machine, which functions in a similar way but allows for
% transitions to various states depending on input; different combinations of
% coins and button presses will unlock different mechanisms in the machine and
% let a specific item fall into the bottom compartment.

% \begin{figure}
%   \begin{tikzpicture}[shorten >=1pt,node distance=5cm,on grid,auto] 
%     \node[text width=2cm, text centered, state] (locked)   {$locked$}; 
%     \node[text width=2cm, text centered, state](unlocked) [right=of locked] {$unlocked$};
%     \node[text centered, node distance=2cm] (coin) [above left=of locked] {$coin$};
%     \path[->] 
%     (coin) edge [] node {} (locked)
%     (locked) edge  [loop left] node {push} ()
%     edge  [bend left=20] node {coin} (unlocked)
%     (unlocked) edge  [bend left=20] node {push} (locked)
%     edge [loop right] node {coin} ();
%   \end{tikzpicture}
%   \caption{Transition diagram for a deterministic finite state automaton modeling a turnstile. Adapted from \citet[762]{koshy2004}.} \label{fig:turnstile}
% \end{figure}

% When a stack is added to this setup, the automaton is essentially provided with
% an infinite memory. Transitions will now be based upon the current state, the
% input symbol and the symbol at the top of the stack. An $\epsilon$-transition
% is possible as well, with $\epsilon$ replacing the input symbol. Thus, any
% transition now consists of the following: the consumption of input, the
% transition to a new state itself, and the replacement of the top of the stack
% by any symbol (it is even possible for the current symbol to remain in place).

% Natural language fits within this paradigm. Given an alphabet $\Sigma$ = \{a,
% b, c, d, \ldots\} or even, in the case of Greek, \{$\alpha$, $\beta$, $\gamma$,
% $\delta$, \ldots\} the initial state and top of the stack will give rise to a
% sequence of characters or strings formed from the alphabet, which may be a
% word, a sentence, an paragraph, \ldots In other words, context-free grammar
% provides a simple yet precise method for mathematically describing and studying
% the rules which govern the construction of natural language from smaller
% blocks; one can parse generated strings (in themselves context-free languages)
% and by induction assemble a grammar. It is also possible to feed a string as
% input to a pushdown automaton applying a context-free grammar; this will show
% whether the string is acceptable by the rules of the grammar or not.

% \begin{figure}
%   \begin{center}
%     \begin{tikzpicture}[shorten >=1pt,node distance=3cm,on grid,auto, label distance=0.4cm] 
%       \node[text centered] (input)   {$input$}; 
%       \node[text centered, align=center, state, rectangle, minimum size=2cm] (fscontrol) [right=of input] {finite\\state\\control};
%       \node[text centered] (acceptreject) [right=of fscontrol] {$accept/reject$};
%       \node[text centered,rectangle split, rectangle split parts=3, draw, label=0:$stack$, minimum size=1.2cm] (stack) [below=of fscontrol] {$top$};
%       \path[->] 
%       (input) edge [] node {} (fscontrol)
%       (fscontrol) edge  [] node {} (acceptreject)
%       edge  [bend left=20] node {} (stack)
%       (stack) edge  [bend left=20] node {} (fscontrol);
%     \end{tikzpicture}
%   \end{center}
%   \caption{General structure of a pushdown automaton. Adapted from \citet[220]{hopcroft2001}.} \label{fig:pushdownautomaton}
% \end{figure}

\chapter{Neural network layers}
% \chapter{Source code}
\label{chp:sourcecode}
For the sake of completeness, we have provided the Python source code
used to create the language model in this appendix. A plain-text copy
is bundled with this document in compact disc format for archival
purposes. Running the model creator without correctly setting up the
directory structure and configuration files will simply yield all
sorts of errors. More extensive documentation can be found in the
README file located in the main project directory. The source may
equally be found on GitHub at \url{https://github.com/sinopeus/thrax}.

\section{Unsupervised training}
\subsection{Trainer}
\label{sec:trainer}
\lstinputlisting[language=Python]{/Users/xavier/projects/thrax/bin/unsupervised-training.py}
\subsection{Unsupervised language model}
\label{sec:langmodel}
\lstinputlisting[language=Python]{/Users/xavier/projects/thrax/bin/unsupervised/model.py}
\subsection{Network parameters}
\label{sec:nnetworkparams}
\lstinputlisting[language=Python]{/Users/xavier/projects/thrax/bin/unsupervised/parameters.py}
\subsection{Network graph}
\label{sec:modelcost}
\lstinputlisting[language=Python]{/Users/xavier/projects/thrax/bin/unsupervised/graph.py}

\section{Supervised training}
\subsection{Trainer}
\label{sec:trainer}
\lstinputlisting[language=Python]{/Users/xavier/projects/thrax/bin/supervised-training.py}
\subsection{Supervised language model}
\label{sec:langmodel}
\lstinputlisting[language=Python]{/Users/xavier/projects/thrax/bin/supervised/nn.py}
\subsection{Network parameters}
\label{sec:nnetworkparams}
\lstinputlisting[language=Python]{/Users/xavier/projects/thrax/bin/supervised/parameters.py}
\subsection{Network graph}
\label{sec:modelcost}
\lstinputlisting[language=Python]{/Users/xavier/projects/thrax/bin/supervised/graph.py}

\section{Tagging}
\subsection{Tagger}
\label{sec:trainer}
\lstinputlisting[language=Python]{/Users/xavier/projects/thrax/bin/tagger.py}
\subsection{Tagging network}
\label{sec:langmodel}
\lstinputlisting[language=Python]{/Users/xavier/projects/thrax/bin/tagging/nn.py}
\subsection{Viterbi trellis}
\label{sec:modelcost}
\lstinputlisting[language=Python]{/Users/xavier/projects/thrax/bin/tagging/viterbi.py}
\subsection{Network graph}
\label{sec:modelcost}
\lstinputlisting[language=Python]{/Users/xavier/projects/thrax/bin/tagging/graph.py}

\section{Data structures \& stream generators}
\subsection{Training state keepers}
\label{sec:trainstate}
\lstinputlisting[language=Python]{/Users/xavier/projects/thrax/bin/state.py}
\subsection{Training example stream}
\label{sec:trainexamplestream}
\lstinputlisting[language=Python]{/Users/xavier/projects/thrax/bin/examples.py}
\subsection{Model validator}
\label{sec:trainexamplestream}
\lstinputlisting[language=Python]{/Users/xavier/projects/thrax/bin/validate.py}
\subsection{Dictionary \& corpus}
\label{sec:dict}
\lstinputlisting[language=Python]{/Users/xavier/projects/thrax/bin/lexicon.py}
\end{document}
